\begin{lemma}[A Gaussian far majorant for $K_{\gamma,j}$, and well-definedness of its hard truncations]
Fix $n \geq 1$ and $j \in \{1,\ldots,n\}$, let $\lambda$ denote Lebesgue measure
on $\mathbb{R}^n$, let $\gamma_n$ be the Gaussian measure of
\noderef{GaussMeasure}, and let $K_{\gamma,j}$ be the normalized Gaussian Riesz
kernel of \noderef{GaussRieszKernel}, so that
$K_{\gamma,j}(x,y) = \sqrt2\,\pi^{-(n+1)/2}\int_0^1\Phi_{x,y}(r)\,dr$ with
\[
  \Phi_{x,y}(r)
    = \frac{y_j - r x_j}{(1-r^2)^{n/2+1}\sqrt{-\log r}}
      \exp\Bigl(-\frac{|rx-y|^2}{1-r^2}\Bigr),
  \qquad r \in (0,1) .
\]
For $\varepsilon > 0$ and $f : \mathbb{R}^n \to \mathbb{R}$ write
\[
  T_{j,\varepsilon}f(x)
    = \int_{|x-y|>\varepsilon} K_{\gamma,j}(x,y)\,f(y)\,dy ,
\]
the hard truncation at radius $\varepsilon$ of the $j$th coordinate, the
integral being against Lebesgue measure in $y$.  Below, ``$f \in
L^1(\gamma_n)$'' means that the \emph{function} $f : \mathbb{R}^n \to
\mathbb{R}$ is $\gamma_n$-integrable; $f$ is a genuine function and not an
equivalence class, because $T_{j,\varepsilon}f(x)$ is formed from its pointwise
values.  Then:
\begin{enumerate}
\item for all $R > 0$ and $\rho > 0$ there is a constant $C > 0$, depending only
on $n$, $j$, $R$ and $\rho$, such that
  \begin{enumerate}
  \item for all $x,y \in \mathbb{R}^n$ with $|x| \leq R$ and
    $|x-y| \geq \rho$, the function $\Phi_{x,y}$ is Lebesgue integrable on
    $(0,1)$ and
    \[
      \bigl|K_{\gamma,j}(x,y)\bigr| \leq C\,e^{-|y|^2} ;
    \]
  \item for every $f \in L^1(\gamma_n)$ and every $x$ with $|x| \leq R$, the
    function $y \mapsto K_{\gamma,j}(x,y)f(y)$ is Lebesgue integrable on
    $\{y : |x-y| > \rho\}$ and
    \[
      \bigl|T_{j,\rho}f(x)\bigr| \leq C\,\|f\|_{L^1(\gamma_n)} ,
      \qquad \|f\|_{L^1(\gamma_n)} = \int_{\mathbb{R}^n}|f|\,d\gamma_n ;
    \]
  \end{enumerate}
\item for every $f \in L^1(\gamma_n)$ and every $\varepsilon > 0$ the function
$x \mapsto T_{j,\varepsilon}f(x)$ is Borel measurable on $\mathbb{R}^n$; and
\item for every $f \in L^1(\gamma_n)$ and every $x \in \mathbb{R}^n$ the
function $\varepsilon \mapsto T_{j,\varepsilon}f(x)$ is continuous on
$(0,\infty)$.
\end{enumerate}

\emph{Why clause (1)(a) carries the integrability assertion, and (1)(b) its
own.}  $K_{\gamma,j}(x,y)$ is defined in \noderef{GaussRieszKernel} as a
Bochner integral over $(0,1)$, with the convention that the value is $0$ when
the integrand is not integrable; and $0$ satisfies every upper bound.  A bare
size bound $|K_{\gamma,j}| \leq Ce^{-|y|^2}$ would therefore be provable
without establishing anything about the kernel, and would be useless to a
consumer that needs the integral to converge.  The same applies one level up:
$T_{j,\rho}f(x)$ is itself a Bochner integral in $y$, so the bound on it is
stated only alongside the assertion that its integrand is integrable.  Both
integrability assertions are part of the content, not hypotheses.  The clauses
(2) and (3) need no such companion: measurability and continuity are not
quantities with a conventional value.

\emph{The majorant carries no polynomial factor in $|y|$.}  The right-hand side
of (1)(a) is $Ce^{-|y|^2}$ with $C$ independent of $y$.  This is what makes
(1)(b) true with $\|f\|_{L^1(\gamma_n)}$ on the right, since
$\|f\|_{L^1(\gamma_n)} = \pi^{-n/2}\int|f(y)|e^{-|y|^2}\,dy$ pairs exactly with
$e^{-|y|^2}$; a residual factor $|y|^k$ would break that pairing.  The constant
$C$ does depend on $R$ and on $\rho$, and no uniformity in $n$, in $R$ or in
$\rho$ is claimed.

\emph{Notation.}  $T_{j,\varepsilon}$ is the $j$th coordinate of the hard
truncation $T_\varepsilon$ displayed in the reference paper.  The paper's
maximal truncation is $\sup_{\varepsilon>0}|T_\varepsilon f(x)|$ with the
\emph{Euclidean vector} norm, which is a different object from the
coordinatewise $\sup_{\varepsilon>0}|T_{j,\varepsilon}f(x)|$; only the
coordinatewise truncations appear here, and the two are reconciled later by
summing over the finitely many coordinates.  Nothing below asserts anything
about the vector object.
\end{lemma}
\begin{proof}
Throughout $n \geq 1$ (the index $j$ ranges over a nonempty set),
$c_n = \sqrt2\,\pi^{-(n+1)/2}$, and $\lambda$ is Lebesgue measure.

\emph{Step 0: the splitting identity.}
For all $x,y \in \mathbb{R}^n$ and $r \in (0,1)$,
\begin{equation*}
  \frac{|rx-y|^2}{1-r^2} = \bigl(|y|^2-|x|^2\bigr) + \frac{|x-ry|^2}{1-r^2} .
  \tag{1}
\end{equation*}
Indeed, multiplying by $1-r^2 > 0$, the claim is
$|rx-y|^2 - (1-r^2)(|y|^2-|x|^2) = |x-ry|^2$, and expanding both sides in the
inner product gives $|x|^2 - 2r\langle x,y\rangle + r^2|y|^2$ for each.
Consequently, writing
\[
  \Psi_{x,y}(r)
    = \frac{y_j-rx_j}{(1-r^2)^{n/2+1}\sqrt{-\log r}}
      \exp\Bigl(-\frac{|x-ry|^2}{1-r^2}\Bigr) ,
\]
we have the pointwise identity
$\Phi_{x,y}(r) = e^{|x|^2}e^{-|y|^2}\,\Psi_{x,y}(r)$ on $(0,1)$.  Since the
factor $e^{|x|^2}e^{-|y|^2}$ is a positive constant in $r$, $\Phi_{x,y}$ is
Lebesgue integrable on $(0,1)$ if and only if $\Psi_{x,y}$ is, and in that case
\begin{equation*}
  K_{\gamma,j}(x,y)
    = c_n\,e^{|x|^2}e^{-|y|^2}\int_0^1\Psi_{x,y}(r)\,dr . \tag{2}
\end{equation*}

\emph{Step 1: a reduction.}
Let $R,\rho > 0$ be given and put $R' = \max(R,1) \geq 1$ and
$\rho' = \min(\rho,1) \in (0,1]$.  Any pair $(x,y)$ with $|x| \leq R$ and
$|x-y| \geq \rho$ also satisfies $|x| \leq R'$ and $|x-y| \geq \rho'$, and
$\|f\|_{L^1(\gamma_n)}$ does not involve $R$ or $\rho$; so a constant that works
for $(R',\rho')$ works for $(R,\rho)$.  We may therefore assume
\[
  R \geq 1, \qquad 0 < \rho \leq 1 .
\]
Fix such $R,\rho$ and fix $x,y$ with $|x| \leq R$ and $|x-y| \geq \rho$.  We
bound $\int_0^1|\Psi_{x,y}|\,d\lambda$ by a constant $C_0 = C_0(n,R,\rho)$ that
does not depend on $x$ or $y$.  We will use throughout that
$|y_j - rx_j| \leq |y| + |x| \leq |y| + R$ for $r \in (0,1)$.

\emph{Step 2: the range $r \in (0,\tfrac12]$.}
Here $1-r^2 \in [\tfrac34,1)$, so $(1-r^2)^{-n/2-1} \leq (4/3)^{n/2+1}$; and
$-\log r \geq \log 2$, so $(-\log r)^{-1/2} \leq (\log 2)^{-1/2}$.  Also
$1-r^2 \leq 1$ gives
$\exp(-|x-ry|^2/(1-r^2)) \leq \exp(-|x-ry|^2)$, and
$|x-ry| \geq r|y| - |x| \geq r|y| - R$, so with $s^+ = \max(s,0)$,
\[
  |\Psi_{x,y}(r)|
    \leq (4/3)^{n/2+1}(\log2)^{-1/2}\,(|y|+R)\,
         \exp\bigl(-((r|y|-R)^+)^2\bigr) .
\]
If $|y| \leq 4R$ then $|y| + R \leq 5R$ and, bounding the exponential by $1$,
\[
  \int_0^{1/2}(|y|+R)\exp\bigl(-((r|y|-R)^+)^2\bigr)\,dr \leq \tfrac52 R .
\]
If $|y| > 4R$ then $|y| > 0$ and the change of variables $t = r|y|$ (a dilation
of $\mathbb{R}$, under which Lebesgue measure scales by $|y|$) gives
\[
  \int_0^{1/2}(|y|+R)\exp\bigl(-((r|y|-R)^+)^2\bigr)\,dr
    = \frac{|y|+R}{|y|}\int_0^{|y|/2}\exp\bigl(-((t-R)^+)^2\bigr)\,dt
    \leq \tfrac54\Bigl(R + \tfrac{\sqrt\pi}{2}\Bigr) ,
\]
using $(|y|+R)/|y| = 1 + R/|y| \leq \tfrac54$ and
$\int_0^\infty\exp(-((t-R)^+)^2)\,dt
 = \int_0^R 1\,dt + \int_R^\infty e^{-(t-R)^2}dt = R + \tfrac{\sqrt\pi}{2}$,
the last integral being the half-line Gaussian integral (shared context: the
scalar Gaussian integral together with the invariance of Lebesgue measure under
the reflection $t \mapsto -t$, which makes the two half-lines contribute
equally).  In either case
\begin{equation*}
  \int_0^{1/2}|\Psi_{x,y}|\,d\lambda
    \leq (4/3)^{n/2+1}(\log2)^{-1/2}\Bigl(\tfrac52R + \tfrac54 R
      + \tfrac{5\sqrt\pi}{8}\Bigr) =: C_1(n,R) . \tag{3}
\end{equation*}

\emph{Step 3: the range $r \in [\tfrac12,1)$.}
Write $\tau = 1-r \in (0,\tfrac12]$.  Then $1-r^2 = \tau(1+r) \in [\tau,2\tau]$,
and $\log r = \log(1-\tau) \leq -\tau$ gives $-\log r \geq \tau$.  Hence
\[
  \frac{1}{(1-r^2)^{n/2+1}\sqrt{-\log r}} \leq \tau^{-n/2-3/2} ,
  \qquad
  \exp\Bigl(-\frac{|x-ry|^2}{1-r^2}\Bigr) \leq \exp\Bigl(-\frac{|x-ry|^2}{2\tau}\Bigr) .
\]
We use repeatedly the elementary identity, valid for $a > 1$ and $b > 0$,
\begin{equation*}
  \int_0^\infty \tau^{-a}e^{-b/\tau}\,d\tau = b^{1-a}\,\Gamma(a-1) < \infty ,
  \tag{4}
\end{equation*}
obtained from the substitution $\tau = b/u$, under which
$d\tau = (b/u^2)\,du$ and the integral becomes
$b^{1-a}\int_0^\infty u^{a-2}e^{-u}\,du$; the latter converges because
$a - 2 > -1$ and because $e^{-u}$ decays at infinity.  We apply $(4)$ with
$a = \tfrac{n+3}{2} > 1$, so $a - 1 = \tfrac{n+1}{2}$.

\emph{Case $|y| \leq 4R$.}  Then $|y_j-rx_j| \leq 5R$.  Since
$x - ry = (x-y) + \tau y$ we get
$|x-ry| \geq |x-y| - \tau|y| \geq \rho - 4R\tau$.  Put
$\tau_0 = \min(\tfrac12,\ \rho/(8R)) > 0$.  For $0 < \tau \leq \tau_0$ we have
$4R\tau \leq \rho/2$, hence $|x-ry| \geq \rho/2$ and
$\exp(-|x-ry|^2/(2\tau)) \leq \exp(-\rho^2/(8\tau))$, so by $(4)$ with
$b = \rho^2/8$,
\[
  \int_{\{0<\tau\leq\tau_0\}} 5R\,\tau^{-n/2-3/2}e^{-\rho^2/(8\tau)}\,d\tau
    \leq 5R\,(\rho^2/8)^{-(n+1)/2}\,\Gamma\bigl(\tfrac{n+1}{2}\bigr) .
\]
For $\tau_0 < \tau \leq \tfrac12$ we bound the exponential by $1$ and
$\tau^{-n/2-3/2} \leq \tau_0^{-n/2-3/2}$, so that part of the integral is at
most $\tfrac52 R\,\tau_0^{-n/2-3/2}$.

\emph{Case $|y| > 4R$.}  Then $|y_j-rx_j| \leq |y| + R \leq \tfrac54|y|$, and
for $r \geq \tfrac12$,
$|x-ry| \geq r|y| - |x| \geq \tfrac12|y| - R > \tfrac12|y| - \tfrac14|y|
 = \tfrac14|y|$, so $\exp(-|x-ry|^2/(2\tau)) \leq \exp(-|y|^2/(32\tau))$.  By
$(4)$ with $b = |y|^2/32$,
\[
  \int_{\{0<\tau\leq 1/2\}}\tfrac54|y|\,\tau^{-n/2-3/2}e^{-|y|^2/(32\tau)}\,d\tau
    \leq \tfrac54|y|\,(|y|^2/32)^{-(n+1)/2}\Gamma\bigl(\tfrac{n+1}{2}\bigr)
    = \tfrac54\,32^{(n+1)/2}\Gamma\bigl(\tfrac{n+1}{2}\bigr)\,|y|^{-n} ,
\]
and $|y| > 4R \geq 4 > 1$ with $n \geq 1$ gives $|y|^{-n} \leq 1$.

Adding the two cases (only one occurs, but the sum dominates either) and
substituting $\tau = 1-r$, which is an isometry of $\mathbb{R}$ composed with a
translation and therefore preserves Lebesgue measure (shared context),
\begin{equation*}
  \int_{1/2}^{1}|\Psi_{x,y}|\,d\lambda \leq C_2(n,R,\rho) , \tag{5}
\end{equation*}
where $C_2 = 5R(\rho^2/8)^{-(n+1)/2}\Gamma(\tfrac{n+1}{2})
 + \tfrac52R\tau_0^{-n/2-3/2} + \tfrac54\,32^{(n+1)/2}\Gamma(\tfrac{n+1}{2})$
depends only on $n$, $R$ and $\rho$.

\emph{Step 4: clause (1)(a).}
By $(3)$ and $(5)$, $\int_0^1|\Psi_{x,y}|\,d\lambda \leq C_0 := C_1 + C_2$,
which is finite and independent of $x$ and $y$ in the range considered.  The
function $\Psi_{x,y}$ is continuous on $(0,1)$, hence Borel measurable there,
so it is Lebesgue integrable on $(0,1)$; by Step 0 so is $\Phi_{x,y}$.  By
$(2)$ and $|x| \leq R$,
\[
  |K_{\gamma,j}(x,y)|
    \leq c_n\,e^{R^2}e^{-|y|^2}\int_0^1|\Psi_{x,y}|\,d\lambda
    \leq c_n\,e^{R^2}C_0\,e^{-|y|^2} .
\]
Set $C_{\mathrm{maj}} = c_n e^{R^2}C_0 > 0$.

\emph{Step 5: joint measurability of the kernel.}
The map
$(x,y,r) \mapsto \Phi_{x,y}(r)$ is continuous on
$\mathbb{R}^n\times\mathbb{R}^n\times(0,1)$: on that domain $1-r^2 \in (0,1]$
and $-\log r \in (0,\infty)$, so every factor is a composition of continuous
functions with nonvanishing denominators.  Hence it is Borel measurable for the
product $\sigma$-algebra.  A partial Bochner integral of a jointly measurable
function is measurable in the remaining parameters, with the convention value
$0$ on the parameters whose slice is not integrable (this is the measurability
half of the Fubini package, and uses only $\sigma$-finiteness of Lebesgue
measure on $(0,1)$, shared context).  Therefore
\begin{equation*}
  (x,y) \longmapsto K_{\gamma,j}(x,y)
  \quad\text{is Borel measurable on } \mathbb{R}^n\times\mathbb{R}^n . \tag{6}
\end{equation*}

\emph{Step 6: clause (1)(b).}
Let $f \in L^1(\gamma_n)$ and let $x$ satisfy $|x| \leq R$.  Being
$\gamma_n$-integrable, $f$ agrees $\gamma_n$-almost everywhere with a Borel
measurable function $\widetilde f$; by clause (4) of
\noderef{GaussLebesgueBallComparison} the measures $\gamma_n$ and $\lambda$ have
the same null sets, so $f = \widetilde f$ also $\lambda$-almost everywhere, and
in particular $f$ is $\lambda$-measurable up to a null set.  By $(6)$ the
function $y \mapsto K_{\gamma,j}(x,y)$ is Borel, so
$y \mapsto K_{\gamma,j}(x,y)f(y)$ is $\lambda$-measurable up to a null set, and
the same is true of its restriction to $\{y : |x-y| > \rho\}$.

On that set $|x-y| \geq \rho$, so clause (1)(a) applies and gives
$|K_{\gamma,j}(x,y)f(y)| \leq C_{\mathrm{maj}}e^{-|y|^2}|f(y)|$.  By
\noderef{GaussMeasure}, $\gamma_n$ is $\lambda$ weighted by
$\pi^{-n/2}e^{-|y|^2}$, so
\begin{equation*}
  \int_{\mathbb{R}^n}e^{-|y|^2}|f(y)|\,dy
    = \pi^{n/2}\int_{\mathbb{R}^n}|f|\,d\gamma_n
    = \pi^{n/2}\|f\|_{L^1(\gamma_n)} < \infty , \tag{7}
\end{equation*}
the finiteness because $f \in L^1(\gamma_n)$.  Hence
$y \mapsto K_{\gamma,j}(x,y)f(y)$ is dominated on $\{|x-y|>\rho\}$ by an
integrable function, so it is Lebesgue integrable there, and
\[
  |T_{j,\rho}f(x)|
    \leq \int_{\{|x-y|>\rho\}}|K_{\gamma,j}(x,y)|\,|f(y)|\,dy
    \leq C_{\mathrm{maj}}\,\pi^{n/2}\,\|f\|_{L^1(\gamma_n)} .
\]
Taking $C = \max\bigl(C_{\mathrm{maj}},\ C_{\mathrm{maj}}\pi^{n/2}\bigr) > 0$
makes both (1)(a) and (1)(b) hold with the single constant $C$, which depends
only on $n$, $j$, $R$ and $\rho$.  This proves clause (1).

\emph{Step 7: clause (2).}
Fix $f \in L^1(\gamma_n)$ and $\varepsilon > 0$, and let $\widetilde f$ be a
Borel function with $f = \widetilde f$ $\lambda$-almost everywhere (Step 6).
For every $x$ the integrands $y \mapsto \mathbf 1_{\{|x-y|>\varepsilon\}}
K_{\gamma,j}(x,y)f(y)$ and the same with $\widetilde f$ agree
$\lambda$-almost everywhere, so they have the same integral; it therefore
suffices to prove the measurability for $\widetilde f$.  The set
$\{(x,y) : |x-y| > \varepsilon\}$ is open in
$\mathbb{R}^n\times\mathbb{R}^n$, hence Borel, so by $(6)$ the function
\[
  (x,y) \longmapsto \mathbf 1_{\{|x-y|>\varepsilon\}}\,K_{\gamma,j}(x,y)\,
  \widetilde f(y)
\]
is Borel on the product.  Applying again the measurability half of the Fubini
package (Step 5), this time integrating out $y$ against the $\sigma$-finite
measure $\lambda$, the function
$x \mapsto \int \mathbf 1_{\{|x-y|>\varepsilon\}}K_{\gamma,j}(x,y)
 \widetilde f(y)\,dy = T_{j,\varepsilon}\widetilde f(x)$ is Borel measurable.

\emph{Step 8: clause (3).}
Fix $f \in L^1(\gamma_n)$ and $x \in \mathbb{R}^n$, and fix
$\varepsilon_* \in (0,\infty)$; we show continuity at $\varepsilon_*$ relative
to $(0,\infty)$.  Put $R_x = \max(|x|,1)$ and $\rho_* = \varepsilon_*/2 > 0$,
and let $C$ be the constant of clause (1) for the pair $(R_x,\rho_*)$.  By
clause (1)(a) and $(7)$ the function
\[
  g(y) = \mathbf 1_{\{|x-y|>\rho_*\}}\,C\,e^{-|y|^2}\,|f(y)|
\]
is Lebesgue integrable, and for every $\varepsilon \geq \rho_*$,
\[
  \bigl|\mathbf 1_{\{|x-y|>\varepsilon\}}K_{\gamma,j}(x,y)f(y)\bigr| \leq g(y)
  \qquad\text{for every } y ,
\]
because $\{|x-y|>\varepsilon\} \subseteq \{|x-y|>\rho_*\}$ and clause (1)(a)
applies on the latter set.

Let $(\varepsilon_k)$ be any sequence in $(0,\infty)$ with
$\varepsilon_k \to \varepsilon_*$; discarding finitely many terms we may assume
$\varepsilon_k \geq \rho_*$ for all $k$.  For every $y$ with
$|x-y| \neq \varepsilon_*$ we have
$\mathbf 1_{\{|x-y|>\varepsilon_k\}} \to \mathbf 1_{\{|x-y|>\varepsilon_*\}}$:
if $|x-y| > \varepsilon_*$ then $|x-y| > \varepsilon_k$ for large $k$, and if
$|x-y| < \varepsilon_*$ then $|x-y| < \varepsilon_k$ for large $k$.  The
exceptional set $\{y : |x-y| = \varepsilon_*\}$ is a sphere in $\mathbb{R}^n$
with $n \geq 1$, hence Lebesgue-null.  So the integrands converge
$\lambda$-almost everywhere and are dominated by the integrable $g$; dominated
convergence (shared context) gives
$T_{j,\varepsilon_k}f(x) \to T_{j,\varepsilon_*}f(x)$.  Since the sequence was
arbitrary and $(0,\infty)$ is a metric space, $\varepsilon \mapsto
T_{j,\varepsilon}f(x)$ is continuous at $\varepsilon_*$ relative to
$(0,\infty)$; as $\varepsilon_*$ was arbitrary, it is continuous on
$(0,\infty)$.
\end{proof}
